%!TEX program = xelatex
\documentclass{xmu}
\begin{document}

% 基础信息

\print % 电子版 / 打印版（打印版将在某些偶数页产生空白页）
% \arabicchapter % 取消注释后，仅使用数字作为章的编号
% \design % 毕业设计 / 毕业论文（取消注释即为毕业设计）
% \minor % 主修 / 辅修（取消注释即为辅修）
\title{中文标题}{English Title}
\author{周淡茗}
\idn{22920212204505}
\college{信息学院}
\subject{你的专业}
\grade{你的年级}
\teacher{校内指导教师\; 职称}
\otherteacher{校外指导教师\; 职称} % 注释则不显示校外指导教师
\pubdate{完成时间}
\keywords{中文关键词}{English keywords}

% 封面、承诺书

\maketitle

% 致谢

\begin{acknowledgement}
    致谢内容。
\end{acknowledgement}

% 中文摘要

\begin{abstract}
    摘要内容。
\end{abstract}

% 英文摘要

\begin{enabstract}
    Abstract Contents.
\end{enabstract}

% 目录

\tableofcontents

% 正文% 正文

\xmuchapter{摘要}{Abstract}
正文内容，脚注\footnote{脚注内容}，引用参考文献\cite{cite1}。

\xmuchapter{绪论}{Introduction}
\xmusection{研究背景与意义}{Research Background and Significance}
正文内容。

\xmusection{国内外研究现状}{Research Status at Home and Abroad}
正文内容。

\xmusection{研究目标与内容}{Research Objectives and Content}
正文内容。

\xmusection{论文组织结构}{Thesis Organization Structure}
正文内容。

\xmuchapter{系统相关技术与理论基础}{System-related Technologies and Theoretical Foundations}
\xmusection{前端技术框架}{Front-end Technology Framework}
正文内容。

\xmusection{后端技术与 API 设计}{Back-end Technology and API Design}
正文内容。

\xmusection{数据通信与状态管理}{Data Communication and State Management}
正文内容。

\xmusection{流程具象化设计}{Process Visualization Design}
\xmusubsection{对话式交互的设计理念}{Design Concept of Dialog-based Interaction}
正文内容。

\xmusubsection{常见的流程具象化实现方式}{Common Process Visualization Implementation Methods}
正文内容。

\xmusection{本章小结}{Summary of This Chapter}
正文内容。

\xmuchapter{系统需求分析}{System Requirements Analysis}
\xmusection{业务需求}{Business Requirements}
正文内容。

\xmusection{功能性需求}{Functional Requirements}
\xmusubsection{增删改查功能}{CRUD Functionality}
正文内容。

\xmusubsection{对话式流程引导}{Dialog-based Process Guidance}
正文内容。

\xmusection{非功能性需求}{Non-functional Requirements}
正文内容。

\xmusection{本章小结}{Summary of This Chapter}
正文内容。

\xmuchapter{系统详细设计}{System Detailed Design}
\xmusection{总体架构设计}{Overall Architecture Design}
正文内容。

\xmusection{核心模块设计}{Core Module Design}
\xmusubsection{增删改查模块}{CRUD Module}
正文内容。

\xmusubsection{对话式流程引导模块}{Dialog-based Process Guidance Module}
正文内容。

%\xmusection{安全设计}{Security Design}
%正文内容。

\xmusection{本章小结}{Summary of This Chapter}
正文内容。

\xmuchapter{系统实现}{System Implementation}
\xmusection{开发环境与工具链}{Development Environment and Toolchain}
正文内容。

\xmusection{核心模块实现}{Core Module Implementation}
\xmusubsection{增删改查功能}{CRUD Functionality}
正文内容。

\xmusubsection{对话式流程引导}{Dialog-based Process Guidance}
正文内容。

\xmusection{本章小结}{Summary of This Chapter}
正文内容。

\xmuchapter{系统测试与验证}{System Testing and Validation}
\xmusection{测试环境}{Testing Environment}
正文内容。

\xmusection{功能测试}{Functional Testing}
正文内容。

\xmusection{安全测试}{Security Testing}
正文内容。

\xmusection{性能测试}{Performance Testing}
正文内容。

\xmusection{本章小结}{Summary of This Chapter}
正文内容。

\xmuchapter{总结与展望}{Conclusion and Outlook}
\xmusection{研究成果总结}{Summary of Research Achievements}
正文内容。

\xmusection{不足与改进方向}{Shortcomings and Improvement Directions}
正文内容。

\xmusection{工作展望}{Future Work Prospects}
正文内容。

\xmuchapter{参考文献}{References}
正文内容。

\xmuchapter{附录}{Appendix}
正文内容。

\xmuchapter{致谢}{Acknowledgments}
正文内容。
% 参考文献

\begin{reference}
    % 如果需要使用 bib 文件导入参考文献，则取消注释下一行
    % \bibliography{references.bib}
    % 已自动按照 GB/T 7714-2005 设置参考文献的引用格式

    % 手动添加参考文献
    \begin{thebibliography}{1} % 括号内数字为参考文献条数
            \bibitem[1]{cite1} 参考文献1
    \end{thebibliography}
\end{reference}

% 附录

\begin{appendix}
    附录内容。
\end{appendix}

% 如果想要将致谢放在最后，可将最前的致谢移动至此

\end{document}